\documentclass[../../main.tex]{subfiles}% 注意这里的文件路径不能用 ./main.tex ,否则用latexmk编译子文件会报错
\graphicspath{{\subfix{./image/}}{\subfix{../../image/}}} % 指定图片引用路径,后续可以直接使用图片文件名
% 如果图片存放在上级目录,则使用 ../../image/ 作为备用路径

% 例如:
% \begin{figure}[H]
% \centering
% \includegraphics[width=0.3\textwidth]{图.png}
% \caption{}
% \label{figure:图}
% \end{figure}
% 注意:上述\label{}一定要放在\caption{}之后,否则引用图片序号会只会显示??.

\begin{document}

\section{其他}

\begin{example}
如果 \( f \in \mathbb{R}[x] \) 是非负的, 证明: 存在 \( \varphi, \psi \in \mathbb{R}[x] \) 使得
\[
f = \varphi^2 + \psi^2.
\]
\end{example}
\begin{proof}
当 \( f \) 为常数多项式则显然, 下面假设 \( \deg f \geqslant 1 \).

设
\[
f(x) = a_0 (x - \alpha_1)^{k_1} (x - \alpha_2)^{k_2} \cdots (x - \alpha_r)^{k_r} g(x),
\]
这里 \( a_0 > 0 \) 且 \( k_i, i = 1,2,\cdots,r \) 都是正偶数且 \( g \) 没有实根.

设
\[
g(x) = (x - \beta_1)(x - \overline{\beta_1})(x - \beta_2)(x - \overline{\beta_2}) \cdots (x - \beta_s)(x - \overline{\beta_s}),
\]
我们有
\[
(x - \beta_1)(x - \beta_2) \cdots (x - \beta_s) = \mu(x) + i\nu(x), \mu, \nu \in \mathbb{R}[x].
\]
注意到
\[
(x - \overline{\beta_1})(x - \overline{\beta_2}) \cdots (x - \overline{\beta_s}) = \mu(x) - i\nu(x),
\]
于是
\[
g(x) = \mu^2(x) + \nu^2(x).
\]

考虑
\[
\varphi(x) \triangleq \sqrt{a_0} \prod_{i=1}^r (x - \alpha_i)^{\frac{k_i}{2}} \mu(x), \psi(x) \triangleq \sqrt{a_0} \prod_{i=1}^r (x - \alpha_i)^{\frac{k_i}{2}} \nu(x),
\]
就有
\[
f = \varphi^2 + \psi^2.
\]
\end{proof}

\begin{example}
设 \( f \) 是 \( \mathbb{C} \) 上无重根的非常数多项式且满足
\[
f(f(x)) = f^n(x) + a_{n-1}f^{n-1}(x) + \cdots + a_1f(x) + a_0,
\]
这里 \( a_0, a_1, \cdots, a_{n-1} \in \mathbb{Z} \). 证明 \( f \in \mathbb{Z}[x] \) 且若 \( a_0, a_1, \cdots, a_{n-1} \) 为奇数, 则 \( f \) 无偶整数根.
\end{example}
\begin{proof}
因为 \( f \) 值域包含无穷多个点, 并且由条件知
\[
f(x), x^n + a_{n-1}x^{n-1} + \cdots + a_1x + a_0,
\]
在无穷多个点上相等, 所以
\[
f(x) = x^n + a_{n-1}x^{n-1} + \cdots + a_1x + a_0 \in \mathbb{Z}[x].
\]
若 \( a_0, a_1, \cdots, a_{n-1} \) 为奇数, 由有理根的性质, 我们知道 \( f \) 的整数根 \( q \) 必须满足 \( q|a_0 \), 从而 \( q \) 为奇数, 这就完成了证明.
\end{proof}

\begin{proposition}
设 \( f,g \in \mathbb{C}[x] \) 是次数大于等于 1 的互素多项式, 证明必然存在唯一的 \( u,v \in \mathbb{C}[x] \) 使得
\[
f(x)u(x) + g(x)v(x) = 1,
\]
且 \( \deg u < \deg g, \deg v < \deg f \).
\end{proposition}
\begin{proof}

存在性:
因为 \( (f,g) = 1 \), 故必然存在 \( k,h \in \mathbb{C}[x] \) 使得 \( fh + gk = 1 \). 若 \( \deg h > \deg g \), 则做带余除法
\[
h = gq + u, \deg u < \deg g.
\]
于是
\[
f(gq + u) + gk = 1 \implies fu + g(fq + k) = 1.
\]
令 \( v = fq + k \), 则由 \( fu + g(fq + k) = 1 \) 和比较两边次数知 \( \deg v < \deg f \). 现在我们证明了存在性.

唯一性:
若还有 \( u_1, v_1 \in \mathbb{C}[x] \) 满足条件, 则
\[
f(u - u_1) = g(v_1 - v).
\]
又 \( f,g \) 互素, 故$
g|(u - u_1), f|(v_1 - v)$,又$\deg (u-u_1)<\deg g,\,\deg(v-v_1)<\deg f$,故$u = u_1, v = v_1$.
这就证明了唯一性.
\end{proof}

\begin{example}
设 \( \mathbb{F} \) 是一个数域, \( f,g \in \mathbb{F}[x], \deg g \geqslant 1 \), 证明: 存在唯一的 \( (f_0,f_1,\cdots,f_r,0,0,\cdots) \in (\mathbb{C}[x])^\infty \),使得
\[
\deg f_i < \deg g \text{或者} f_i = 0, 0 \leqslant i \leqslant r
\]
且
\[
f(x) = f_0(x) + f_1(x)g(x) + f_2(x)g^2(x) + \cdots + f_r(x)g^r(x).
\]
\end{example}
\begin{remark}
\( (\mathbb{C}[x])^\infty \) 表示 \( \mathbb{C}[x] \) 的可数笛卡尔积. 为书写方便, 我们可以记 \( f_i = 0, i \geqslant r + 1 \).
\end{remark}
\begin{proof}
为方便, 对 \( p \in \mathbb{C}[x] \), 我们约定 \( \deg p = -\infty \Leftrightarrow p = 0 \).

存在性:
由带余除法, 存在 \( r \in \mathbb{N}_0 \) 使得 \( \deg q_r < \deg g \)(否则,$g$就能做无穷次带余除法,但$g$的次数有限) 且
\[
f = q_1g + f_0, \deg f_0 < \deg g, q_1 = q_2g + f_1, \deg f_1 < \deg g
\]
\[
q_2 = q_3g + f_2, \deg f_2 < \deg g, q_3 = q_4g + f_3, \deg f_3 < \deg g
\]
\[
\vdots
\]
\[
q_{r-1} = q_r g + f_{r-1}, \deg f_{r-1} < \deg g,
\]
取 \( f_r = q_r \), 我们得
\[
f = f_0 + (q_2g + f_1)g = f_0 + f_1g + (q_3g + f_2)g^2 = \cdots
\]
\[
= f_0 + f_1g + \cdots + f_{r-1}g^{r-1} + q_r g^r = f_0 + f_1g + \cdots + f_r g^r.
\]

唯一性:
若有两种不同的表示
\[
f = \sum_{i=0}^\infty f_i g^i = \sum_{i=0}^\infty h_i g^i, \deg h_i, \deg f_i < \deg g, i \in \mathbb{N}_0,
\]
我们有
\[
\sum_{i=0}^\infty (f_i - h_i) g^i = 0.
\]
设上式使得 \( f_i - h_i \neq 0 \) 的最大的 \( i \) 为 \( k \), 则 \( \deg [(f_k - h_k) g^k] \geqslant k \deg g \) 以及
\[
\deg [(f_i - h_i) g^i] = \deg (f_i - h_i) + i \deg g < k \deg g, 0 \leqslant i \leqslant k,
\]
故 \( \sum_{i=0}^\infty (f_i - h_i) g^i \) 不可能是 0 多项式(最高次项消不掉), 矛盾! 这就证明了唯一性.
\end{proof}

\begin{example}
设 \( p,q \in \mathbb{R}[x] \) 满足
\[
p(x^2 + x + 1) = q(x^2 - x + 1),
\]
证明 \( p = q \) 为常数.
\end{example}
\begin{note}
利用等式关系反复迭代来得到无穷个根.
\end{note}
\begin{proof}
事实上
\[
(2x + 1)p'(x^2 + x + 1) = (2x - 1)q'(x^2 - x + 1), \forall x \in \mathbb{R},
\]
那么 \( p'\left( \left( \frac{1}{2} \right)^2 + \frac{1}{2} + 1 \right) = 0 \), 因此由(对称轴是$x=-\frac{1}{2}$)Vieta定理知
\[
\left( \frac{1}{2} \right)^2 + \frac{1}{2} + 1 = \left( -\frac{3}{2} \right)^2 - \frac{3}{2} + 1.
\]
从而 \( p'\left( \left( -\frac{3}{2} \right)^2 - \frac{3}{2} + 1 \right) = 0 \), 类似的从而 \( q'\left( \left( -\frac{3}{2} \right)^2 + \frac{3}{2} + 1 \right) = 0 \), (对称轴是$x=\frac{1}{2}$)从而 \( q'\left( \left( \frac{5}{2} \right)^2 - \frac{5}{2} + 1 \right) = 0 \), 从而 $p'$ $\left( \left( \frac{5}{2} \right)^2 + \frac{5}{2} + 1 \right) = 0$, 依次下去我们得到 \( p',q' \) 有无穷多个根, 因此都为 0, 这样结合 \( p(1) = q(1) \) 我们就证明了 \( p = q \) 为常数.
\end{proof}

\begin{example}
计算全部 \( f \in \mathbb{C}[x] \) 使得
\begin{align}
f(x^2) = f(x)f(x+1), \forall x \in \mathbb{C}. \label{23.173}
\end{align}
\end{example}
\begin{proof}
当 \( f = c \in \mathbb{C} \), 我们有 \( c = c^2 \), 故 \( f = 0 \) 或者 \( f = 1 \) 为所求.

当 \( f \) 不为常数, 设 \( f(x_1) = 0, x_1 \in \mathbb{C} \), 我们由\(\eqref{23.173}\)知
\[
f(x_1^2) = 0.
\]
反复运用\(\eqref{23.173}\)得 \( x_1^{2^n}, n \in \mathbb{N}_0 \) 都是 \( f \) 零点. 又非0多项式零点有限, 故设 \( x_1^{2^n} = x_1^{2^m}, n > m \geqslant 0 \), 则 \( x_1^{2^n - 2^m} = 1 \)或 \( x_1 = 0 \). 当前者发生, 我们有 \( |x_1| = 1 \). 从而 \( f \) 的所有根都在单位圆周上或者是 0.

设 \( f(x_1) = 0, x_1 \in \mathbb{C} \), 由\(\eqref{23.173}\)知 
\begin{align*}
f((x-1)^2)=f(x-1)f(x).
\end{align*}
故\( f\left( (x_1 - 1)^2 \right) = 0 \). 同理可得\( |x_1 - 1| = 1 \)或$x_1-1=0$.因此由$|x_1|=1,|x_1-1|=1,x_1=0\text{或}1$可知$f$只可能有$x_1=0,1,e^{\pm \frac{\pi}{3}i} $这四个零点.但是若$x_1=e^{\pm \frac{\pi}{3}i}$是$f$的零点,则由\eqref{23.173}知$x_1^2=e^{\pm \frac{2\pi}{3}i}$也是$f$的零点,矛盾!
从而 \( f \) 的零点只有 0 或者 1. 设 \( f(x) = cx^n(x - 1)^m \) 代入原方程得
\begin{align*}
f(x^2) &= cx^{2n}(x^2 - 1)^m = cx^{2n}(x - 1)^m(x + 1)^m \\
&= f(x)f(x+1) = c^2x^{n+m}(x - 1)^m(x + 1)^n \\
&\implies f(x) = x^n(x - 1)^n, n \in \mathbb{N}.
\end{align*}
这就完成了证明.
\end{proof}

\begin{example}
设$ p \in \mathbb{Z}[x] $使得$ p $在7个不同整数点取值为7, 证明$ p $没有整数根.
\end{example}
\begin{proof}
由条件可设
\begin{align*}
p\left( x \right) -7=\left( x-x_1 \right) \left( x-x_2 \right) \cdots \left( x-x_7 \right) q\left( x \right),
\end{align*}
其中$x_1,x_2,\cdots,x_7$为互不相同的整数，$q\in \mathbb{Z}[x]$.假设$a$为$p\left( x \right)$的整数根，则
\begin{align*}
-7=\left( a-x_1 \right) \left( a-x_2 \right) \cdots \left( a-x_7 \right) q\left( a \right),
\end{align*}
其中$a-x_i$为互不相同的整数.故$\left( a-x_1 \right),\left( a-x_2 \right),\cdots,\left( a-x_7 \right)$是$-7$的$7$个不同因子.这与$-7$只有因子$1,-1,7,-7$这四个不同因子矛盾!
\end{proof}

\begin{example}
设 $\alpha$ 是 $x^4+x^3+x^2+x+1=0$ 的根, $\varphi: \mathbb{Q}[x] \to \mathbb{C}$ 是线性变换, 满足: $\varphi(f(x))=f(\alpha),$ \(\forall f \in \mathbb{Q}[x].\) 证明: $(2+\alpha)^{-1} \in \operatorname{Im} \varphi.$
\end{example}
\begin{proof}
因为 $f(-2)=16-8+4-2+1=11\ne 0,$ 所以 \(x+2\) 与 \(f(x)\)在$\mathbb{Q}$上互素. 从而存在 $u(x), v(x) \in \mathbb{Q}[x],$ 使得
\begin{align*}
u(x)f(x)+v(x)(x+2)=1.
\end{align*}
将 \(x=\alpha\) 代入得 $u(\alpha)\cdot 0+v(\alpha)\cdot(\alpha+2)=1,$ 从而 $v(\alpha)=(\alpha+2)^{-1}=(2+\alpha)^{-1}.$ 由于 $\varphi(v(x))=v(\alpha),$ 故 $(2+\alpha)^{-1}\in \operatorname{Im} \varphi.$\end{proof}

\begin{example}
求多项式 \(f(x),g(x)\), 使得 \(x^{m}f(x)+(1-x)^{n}g(x)=1\).
\end{example}
\begin{solution}

{\color{blue}解法一:} 
首先, 对原式在模 \((1-x)^n\) 下进行运算, 得到
\begin{align}\label{eq:qqqqqqqAA-1}
x^{m}f(x)\equiv 1\pmod{(1-x)^n}.
\end{align}
令 \(t=1-x\), 则 \(x=1-t\), 代入上式得
\begin{align*}
(1-t)^{m}f(1-t)\equiv 1\pmod{t^n}.
\end{align*}
由此可得
\begin{align*}
f(1-t)\equiv (1-t)^{-m}\pmod{t^n}.
\end{align*}
利用Taylor展开公式得
\begin{align*}
(1-t)^{-m}=\sum\limits_{i=0}^{\infty}\binom{m+i-1}{i}t^i.
\end{align*}
上式在模 \((1-x)^n\) 下进行运算得
\begin{align}
f(x)=\sum\limits_{i=0}^{n-1}\binom{m+i-1}{i}(1-x)^i. \label{eq:standard_f}
\end{align}
同理, 对原式在模 \(x^m\) 下进行运算, 得到
\begin{align}\label{eq:qqqqqqqAA-2}
(1-x)^n g(x)\equiv 1\pmod{x^m}.
\end{align}
因此
\begin{align*}
g(x)\equiv (1-x)^{-n}\pmod{x^m}.
\end{align*}
再次利用利用Taylor展开公式并在模 \(x^m\)下进行运算得
\begin{align}
g(x)=\sum\limits_{j=0}^{m-1}\binom{n+j-1}{j}x^j. \label{eq:standard_g}
\end{align}
由\eqref{eq:qqqqqqqAA-1}式可知 $x^{m}f(x)-1$ 能被 $(1-x)^{n}$ 整除, 由\eqref{eq:qqqqqqqAA-2}式可知 $(1-x)^{n}g(x)-1$ 能被 $x^{m}$ 整除.
令
\begin{align*}
E(x)=x^{m}f(x)+(1-x)^{n}g(x)-1,
\end{align*}
则有 $E(x)\equiv 0 \pmod{(1-x)^n}$ 且 $E(x)\equiv 0 \pmod{x^m}$.
由于 $x^m$ 与 $(1-x)^n$ 互素, 故由\subcref{proposition:互素多项式和最大公因式的基本性质}{proposition:互素多项式和最大公因式的基本性质-1}知$x^m(1-x)^n \mid E(x)$.
另一方面, 因 $\deg f \leqslant n-1$, $\deg g \leqslant m-1$, 所以 $\deg E \leqslant m+n-1$.
而 $x^m(1-x)^n$ 的次数为 $m+n$,因此 $E(x) \equiv 0$.
即 $x^{m}f(x)+(1-x)^{n}g(x)=1$.

综上可知,\eqref{eq:standard_f}式和 \eqref{eq:standard_g}式就是原式满足\(\deg f \leqslant n-1\) 且 \(\deg g \leqslant m-1\) 的解 .

{\color{blue}解法二:} 
注意到因式分解公式
\begin{align*}
1-x^m=(1-x)(1+x+x^2+\cdots+x^{m-1}).
\end{align*}
对两边同时取 \(n\) 次方, 可得
\begin{align}
(1-x^m)^n=(1-x)^n(1+x+x^2+\cdots+x^{m-1})^n. \label{eq:identity}
\end{align}
另一方面, 利用二项式定理展开 \((1-x^m)^n\) 得
\begin{align*}
(1-x^m)^n =\sum\limits_{k=0}^{n}\binom{n}{k}(-1)^k x^{km} =1-x^{m}\sum\limits_{k=1}^{n}(-1)^{k-1}\binom{n}{k}x^{(k-1)m}.
\end{align*}
将其代入等式 \eqref{eq:identity}, 便得到
\begin{align*}
1=x^{m}\sum\limits_{k=1}^{n}(-1)^{k-1}\binom{n}{k}x^{(k-1)m}+(1-x)^n(1+x+x^2+\cdots+x^{m-1})^n.
\end{align*}
因此可以直接构造一组特解:
\begin{align*}
f(x)  =\sum\limits_{k=1}^{n}(-1)^{k-1}\binom{n}{k}x^{(k-1)m}, \quad
g(x)  =(1+x+x^2+\cdots+x^{m-1})^n.
\end{align*}
\end{solution}

\begin{example}
设 $\phi(x)$ 是有理数域上的一个不可约多项式，$x_1,x_2,\cdots,x_s$ 是 $\phi(x)$ 在复数域上的根，$f(x)$ 是任一个有理系数多项式，使 $f(x)$ 不能被 $\phi(x)$ 整除。证明：存在有理系数多项式 $h(x)$，使
\begin{align*}
\frac{1}{f(x_i)}=h(x_i),\qquad i=1,2,\cdots,s.
\end{align*}
\end{example}
\begin{proof}

{\color{blue}证法一:}
因为 $\phi(x)$ 在 $\mathbb{Q}[x]$ 中不可约，且 $f(x)$ 不能被 $\phi(x)$ 整除，所以 $\gcd(f(x),\phi(x))=1$。于是由\hyperref[theorem:最大公因式的必要条件]{裴蜀等式}，存在 $u(x),v(x)\in\mathbb{Q}[x]$，使得
\begin{align*}
u(x)f(x)+v(x)\phi(x)=1.
\end{align*}
将 $x=x_i$ 代入，得
\begin{align*}
u(x_i)f(x_i)+v(x_i)\phi(x_i)=1. 
\end{align*}
又 $\phi(x_i)=0$，所以
\begin{align*}
u(x_i)f(x_i)=1.
\end{align*}
因 $u(x_i)f(x_i)=1$，故 $f(x_i)\neq0$，且
\begin{align*}
\frac{1}{f(x_i)}=u(x_i),\quad i=1,\cdots,s. 
\end{align*}
取 $h(x)=u(x)$，则 $h(x)\in\mathbb{Q}[x]$，且满足
\begin{align*}
h(x_i)=\frac{1}{f(x_i)},\quad i=1,\cdots,s. 
\end{align*}

{\color{blue}证法二:}
设 $\deg\phi=s$。由于有理数域的特征为零，不可约多项式 $\phi(x)$ 没有重根，因此它在复数域中的根 $x_1,\cdots,x_s$ 两两不同。

首先证明对每个 $i$ 都有 $f(x_i)\neq0$。若存在某个 $i$ 使 $f(x_i)=0$，则 $x_i$ 同时是 $\phi(x)$ 与 $f(x)$ 的根。因为 $\phi(x)$ 是 $x_i$ 在 $\mathbb{Q}$ 上的不可约多项式，故必有 $\phi(x)\mid f(x)$，这与题设矛盾!因此 $f(x_i)\neq0$，$i=1,\cdots,s$。

对插值节点 $x_1,\cdots,x_s$ 及 $(x_i,\frac{1}{f(x_i)})$ 应用\hyperref[Numerical Analysis-theorem:Lagrange插值多项式111]{Lagrange插值定理}。存在唯一的次数小于 $s$ 的复系数多项式
\begin{align}
h(x)=\sum\limits_{i=1}^{s}\frac{1}{f(x_i)}\prod\limits_{\substack{1\leqslant j\leqslant s\\j\neq i}}\frac{x-x_j}{x_i-x_j}, \label{lagrange}
\end{align}
则显然有$h(x_i)=\frac{1}{f(x_i)}$，$\,i=1,\cdots,s$。

下面证明 $h(x)\in\mathbb{Q}[x]$。将\eqref{lagrange}式写成
\begin{align}
h(x)=c_0+c_1x+\cdots+c_{s-1}x^{s-1}. \label{coef}
\end{align}
由\eqref{lagrange}式可知，每个 $c_k$ 都是根 $x_1,\cdots,x_s$ 的有理函数。任意交换这些根，只会改变\eqref{lagrange}式中各项的排列，而不会改变整个多项式 $h(x)$,于是也不会改变$h(x)$的各个系数$c_k$ ,所以每个 $c_k$ 都是 $x_1,\cdots,x_s$ 的对称有理函数。

再设
\begin{align*}
\phi(x)=a_sx^s+a_{s-1}x^{s-1}+\cdots+a_0,\quad a_i\in\mathbb{Q}.
\end{align*}
由\hyperref[theorem:Vieta定理]{Vieta定理}知
\begin{align*}
e_r(x_1,\cdots,x_s)=(-1)^r\frac{a_{s-r}}{a_s}\in\mathbb{Q},\quad r=1,2,\cdots,s.
\end{align*}
其中$e_r$表示$x_1,\cdots,x_s$的初等对称多项式.
根据\hyperref[Abstract Algebra-theorem:抽象代数--定理2.3.1--对称多项式基本定理]{对称多项式基本定理}，每个 $c_k$ 都可以表示为$e_1,\cdots,e_s$ 的有理函数。从而 $c_k\in\mathbb{Q}$，$k=0,1,\cdots,s-1$。因此 $h(x)\in\mathbb{Q}[x]$，且
\begin{align*}
\frac{1}{f(x_i)}=h(x_i),\quad i=1,2,\cdots,s. 
\end{align*}
\end{proof}

\begin{example}
证明: 数域 \(F\) 上的任意一个 \(n\) 元多项式都可以表示成一次齐次多项式方幂的线性组合.
\end{example}
\begin{proof}
对 \(n\) 作归纳法. 当 \(n=2\) 时, \(f(x_1, x_2) = \sum\limits_{n_1, n_2} a_{n_1, n_2} x_1^{n_1} x_2^{n_2}\), 只需对其中的任一单项式 \(x^r y^s\) (无需考虑系数), 证明结论成立. 为此, 可设
\begin{align}
x^r y^s = k_0(x+y)^{r+s} + k_1(x+2y)^{r+s} + \cdots + k_{r+s}(x+(r+s)y)^{r+s} \label{eq:setup}
\end{align}
将等式右边展开，比较 \(y^j\) (\(j=0,1,\cdots,r+s\)) 的系数，可知 \eqref{eq:setup} 式成立当且仅当以下线性方程组有解：
\begin{align*}
\begin{cases}
k_0+k_1+\cdots +k_{r+s}=0\\
k_0+2k_1+\cdots +\left( r+s \right) k_{r+s}=0\\
\cdots \cdots \cdots \cdots\\
k_0+2^sk_1+\cdots +\left( r+s \right) ^sk_{r+s}=\frac{1}{\mathrm{C}_{r+s}^{s}}\\
\cdots \cdots \cdots \cdots\\
k_0+2^{r+s}k_1+\cdots +\left( r+s \right) ^{r+s}k_{r+s}^{s}=0\\
\end{cases}
\end{align*}
显然，该线性方程组的系数矩阵为Vandermode矩阵，由于 \(1, 2, \cdots, r+s\) 互不相同，其系数行列式不为零. 根据 Cramer 法则，上述方程组有唯一解. 因此，\(n=2\) 时结论成立.

假设对变量个数 \(\leqslant n-1\) 的多项式结论成立. 对于 \(n\) 元多项式 \(f(x_1,x_2,\cdots,x_n)\)，我们有
\begin{align}
f(x_1,x_2,\cdots,x_n) = \sum\limits_{s} f_s(x_1,x_2,\cdots,x_{n-1}) x_n^s, \label{eq:induction_step}
\end{align}
其中 \(f_s(x_1,x_2,\cdots,x_{n-1})\) 是变量个数不超过 \(n-1\) 的多项式. 根据归纳假设，可用一次齐次多项式幂的线性组合表示为
\begin{align*}
f_s(x_1,x_2,\cdots,x_{n-1}) = \sum\limits_{r} a_{rs} (c_{s1} x_1 + c_{s2} x_2 + \cdots + c_{s,n-1} x_{n-1})^r. 
\end{align*}
代入 \eqref{eq:induction_step} 式，得
\begin{align}
f(x_1,x_2,\cdots,x_n) = \sum\limits_{s} \sum\limits_{r} a_{rs} (c_{s1} x_1 + c_{s2} x_2 + \cdots + c_{s,n-1} x_{n-1})^r x_n^s. \label{eq:combined}
\end{align}
将上式右端的每一个括号视为一个变量，记为 \(z_i\)，对二元单项式 \(z_i^r x_n^s\) 利用已证明的事实，有
\begin{align*}
z_i^r x_n^s = \sum\limits_{i=0}^{r+s} k_i (z_i + (i+1) x_n)^{r+s}. 
\end{align*}
代入 \eqref{eq:combined} 式右端，即为 \(x_1, x_2, \cdots, x_n\) 的一次齐次多项式幂的线性组合. 因此，命题得证.
\end{proof}



\end{document}